\src{108}
đối tích của hai vật không bao giờ tồn tại, và phép đẩy ra của hai vật $A,B$
dưới một vật thứ ba $C$ nói chung không tồn tại (chẳng hạn, lấy
$C=\Delta_0$, $A=B=\Delta_1$, với hai ánh xạ cấu trúc $u:C\to A$ và
$v:C\to B$ bằng nhau). Tuy nhiên, nó tồn tại trong một số trường hợp; chẳng
hạn
\[
  C=\Delta_0,\qquad A=\Delta_m,\qquad B=\Delta_n,\qquad
  u(0)=m,\qquad v(0)=0,
\]
và trong trường hợp này
\[
  \Delta_m\amalg_{\Delta_0}\Delta_n=\Delta_{m+n}.
\]

Một \emph{vật đơn hình} (tương ứng, một \emph{vật bán đơn hình}) trong một
phạm trù $\mathcal C$ theo định nghĩa là một hàm tử phản biến $K$ từ
$\mathcal S$ (tương ứng từ $\mathcal S'$) vào $\mathcal C$. Do đó, một vật
đơn hình xác định một vật bán đơn hình bằng phép hạn chế, nhưng khái niệm thứ
nhất khác khái niệm thứ hai về căn bản bởi sự hiện diện của \emph{các phép
toán đối xứng} trên các $K_n=K(\Delta_n)$, tương ứng với ảnh qua hàm tử $K$
của các phần tử thuộc nhóm đối xứng trên $n+1$ phần tử (được xem là nhóm các
tự đẳng cấu của $\Delta_n$ trong $\mathcal S$).

Với mọi $n$, gọi $\Delta'_n$ (tương ứng $\Delta''_n$) là phạm trù hữu hạn có
tập các vật là $\Delta_n$, được xác định bởi quan hệ thứ tự hỗn loạn (tương
ứng, thứ tự toàn phần tự nhiên) trên $\Delta_n$ (tập các mũi tên là đồ thị của
quan hệ thứ tự ấy). Hiển nhiên $\Delta'_n$ (tương ứng $\Delta''_n$) phụ thuộc
theo hàm tử vào vật $\Delta_n$ của $\mathcal S$ (tương ứng $\mathcal S'$).
Nếu $Z$ là một phạm trù, thì $\operatorname{Hom}(\Delta'_n,Z)$ (tương ứng
$\operatorname{Hom}(\Delta''_n,Z)$), khi $\Delta_n$ biến thiên, là một hàm tử
từ phạm trù $\mathcal S$ (tương ứng $\mathcal S'$) vào phạm trù các tập hợp,
tức một tập đơn hình (tương ứng, một tập bán đơn hình), được gọi là
\emph{liên kết với phạm trù $Z$}; ký hiệu chúng lần lượt là $Z'$ và $Z''$.
Hơn nữa, có một đồng cấu tự nhiên hiển nhiên từ tập bán đơn hình liên kết
với $Z'$ vào
$Z''$, và đồng cấu này là một đẳng cấu khi và chỉ khi $Z$ là một groupoid. Với
các quy ước này:

\result{MỆNH ĐỀ}{4.1} Hàm tử $Z\longmapsto Z''$ từ phạm trù các phạm trù vào
phạm trù các tập bán đơn hình là đầy và trung thành, và xác định một tương
đương giữa phạm trù các \emph{phạm trù} với phạm trù các tập bán đơn hình,
tức các hàm tử phản biến $K$ từ $\mathcal S'$ vào $(\mathrm{Set})$, biến các
phép đẩy ra $A\amalg_C B$ thuộc kiểu được chỉ ra ở trên thành các tích thớ của
các tập hợp. Tương tự, hàm tử $Z\longmapsto Z'$ từ phạm trù các groupoid vào
phạm trù các tập đơn hình là đầy và trung thành, và xác định một tương
đương giữa phạm trù các \emph{groupoid} với phạm trù các
